\section{Khởi động và Kích hoạt Phần mềm}
    Sau khi hoàn tất quá trình lắp ráp phần cứng, hệ thống gateway (với kiến trúc Dual-MCU gồm ESP32-S3 Master/WAN-MCU và Slave/LAN-MCU) được đưa vào vận hành theo chu trình: Khởi động – Nạp cấu hình – Kết nối – Vận hành dịch vụ.

    \subsection{Khởi động hệ thống (Boot \& RTOS Init)}

    Khi được cấp nguồn, WAN-MCU đóng vai trò điều phối chính (Application Master). Vi điều khiển này thực hiện khởi tạo các thành phần phần cứng nền tảng (Drivers, BSP) và khởi chạy các tác vụ (tasks) trên hệ điều hành FreeRTOS theo mô hình handler chạy song song. Các khối chức năng chính được kích hoạt bao gồm:

    \begin{itemize}
        \item \textbf{Main Control:} Điều phối trạng thái hoạt động của toàn hệ thống.
        \item \textbf{Internet Communication Handler:} Quản lý các kết nối mạng (Ethernet, Wi-Fi, LTE/NB-IoT).
        \item \textbf{Server Communication Handler:} Quản lý giao thức tầng ứng dụng (MQTT, HTTP, CoAP).
        \item \textbf{MCU LAN Communication Handler:} Thiết lập kênh giao tiếp liên vi điều khiển (Inter-MCU) với LAN-MCU.
        \item \textbf{FOTA/Bridge Bootloader:} Sẵn sàng cho quy trình cập nhật firmware từ xa.
    \end{itemize}

    \subsection{Nạp và áp dụng cấu hình vận hành (Load/Apply Configuration)}

    Hệ thống truy xuất toàn bộ tham số cấu hình được lưu trữ trong bộ nhớ. Gateway sẽ đọc và áp dụng các cấu hình này cho các khối Internet/Cloud và đặc biệt là tải các driver tương ứng cho các module phần cứng đang được lắp đặt. Cơ chế này đảm bảo tính \textbf{modularity} và \textbf{configurability}, cho phép hệ thống chỉ khởi chạy các tài nguyên cần thiết theo nhu cầu triển khai thực tế.

    \subsection{Thiết lập kết nối WAN và đồng bộ thời gian}

    Dựa trên cấu hình đã tải, gateway tiến hành thiết lập kết nối Internet ưu tiên (Wi-Fi hoặc LTE).

    \begin{itemize}
        \item \textbf{Đối với kết nối Wi-Fi,} hệ thống hỗ trợ cả hai chế độ bảo mật là Personal và Enterprise.
        \item Ngay sau khi kết nối mạng thành công, gateway thực hiện giao thức SNTP để đồng bộ thời gian thực và cập nhật cho RTC, đảm bảo tính chính xác của nhãn thời gian (timestamp) cho dữ liệu log và telemetry.
    \end{itemize}

    \subsection{Kết nối Cloud và vòng lặp vận hành dữ liệu}

    Khi kết nối WAN ổn định, gateway khởi tạo kết nối đến Cloud Server (thông qua giao thức MQTT). Lúc này, hệ thống đi vào vòng lặp vận hành chính:

    \begin{itemize}
        \item \textbf{Luồng Uplink:} Dữ liệu thu thập từ LAN-MCU được truyền sang WAN-MCU, đóng gói (packetizing) và đẩy vào hàng đợi (queue) để gửi lên broker/server theo cấu hình.
        \item \textbf{Luồng Downlink:} Gateway duy trì kết nối để nhận lệnh điều khiển từ server và phân luồng xuống LAN-MCU hoặc xử lý tại chỗ.
        \item \textbf{Cơ chế bảo vệ:} Hệ thống duy trì cơ chế tự động kết nối lại (auto-reconnect) khi mất mạng để đảm bảo tính sẵn sàng cao.
    \end{itemize}

    \subsection{Kích hoạt cấu hình tại chỗ (Local Configuration)}

    Trong trường hợp cần cấu hình thủ công bởi kỹ thuật viên, gateway hỗ trợ chế độ cấu hình qua cổng Programmer USB-C (giao tiếp USB/UART).

    \begin{itemize}
        \item \textbf{Scan/Read Config:} PC gửi yêu cầu đọc, gateway trả về các cặp key–value hiện hành (có cơ chế che thông tin nhạy cảm).
        \item \textbf{Write Config:} PC gửi lệnh cấu hình (prefix “CF…”), gateway đưa vào hàng đợi xử lý $\to$ kiểm tra tính hợp lệ (validate) $\to$ lưu vào NVS $\to$ áp dụng ngay (runtime apply). Nếu cấu hình thuộc về miền LAN, WAN-MCU sẽ chuyển tiếp lệnh xuống LAN-MCU để xử lý đồng bộ.
    \end{itemize}

    \subsection{Cập nhật Firmware (FOTA)}

    Gateway hỗ trợ quy trình cập nhật firmware an toàn (Safety FOTA) nhằm giảm thiểu rủi ro hỏng thiết bị (brick). Quy trình bao gồm kiểm tra tính toàn vẹn (checksum), thực hiện cập nhật tuần tự cho từng MCU (Master và Slave) và cơ chế không ghi đè lên firmware đang chạy cho đến khi gói cập nhật được xác thực thành công.
